% 三版共用的扉页：书名、英文书名、作者、配图，以及底部的仓库地址与 Last commit；设了版本标记时悬挂在书名右侧。
% 由各版布局在 shared/fonts.tex、shared/colors.tex 之后 \input，再用 \TitlePageSetup 给出本版的版面参数。
% 书名、作者、英文书名、配图路径、版本标记与仓库地址由入口提供；路径均相对于项目根目录。
% 绝对定位依赖 TikZ 的 remember picture，至少编译两轮（根目录 run.sh 已保证）。

\newfontfamily\BookTitlePageMediumLatin{PingFang SC Medium}
\newCJKfontfamily\BookTitlePageMediumCJK{PingFang SC Medium}
% 提交信息的眉题：小号大写，加宽字距。
\newfontfamily\BookTitlePageEyebrow{PingFang SC Medium}[LetterSpace=14]
% 提交说明原样输出：此字体族不启用 TeX 连字，-- 等不会被替换。
\newfontfamily\BookTitlePagePlainLatin{PingFang SC}

% 版面参数：以纸张中心为原点，书名、英文书名、作者的基线，配图区中心，底部信息区底边的纵向位置，
% 以及配图区的宽和高。配图在配图区内等比缩放后居中；标题区与底部信息区的位置不随图片比例变化。
% 各版布局须给出全部 7 项，缺少时编译报错；数值按原文交给 TikZ 与 \setlength，不预先换算。
\newlength{\BookTitlePageImageWidth}
\newlength{\BookTitlePageImageHeight}
\pgfkeys{
  /book/titlepage/.is family,
  /book/titlepage/title/.store in=\BookTitlePageTitleY,
  /book/titlepage/english/.store in=\BookTitlePageEnglishY,
  /book/titlepage/author/.store in=\BookTitlePageAuthorY,
  /book/titlepage/image/.store in=\BookTitlePageImageY,
  /book/titlepage/stamp/.store in=\BookTitlePageStampY,
  /book/titlepage/image-width/.code={\setlength{\BookTitlePageImageWidth}{#1}},
  /book/titlepage/image-height/.code={\setlength{\BookTitlePageImageHeight}{#1}},
}
\newcommand{\TitlePageSetup}[1]{\pgfkeys{/book/titlepage/.cd,#1}}

\makeatletter
% 元信息：英文书名、配图路径与版本标记；不写或参数为空时省略对应内容。
% 版本标记只写文字（如“代码版”），全角括号与样式由扉页添加。
\newcommand*{\EnglishTitle}[1]{\gdef\BookTitlePage@EnglishTitle{#1}}
\newcommand*{\TitlePageImage}[1]{\gdef\BookTitlePage@Image{#1}}
\newcommand*{\EditionLabel}[1]{\gdef\BookTitlePage@Edition{#1}}
\let\BookTitlePage@EnglishTitle\@empty
\let\BookTitlePage@Image\@empty
\let\BookTitlePage@Edition\@empty
% PDF 文档属性同步入口的书名与作者；有版本标记时标题附上，如“算法竞赛小记（代码版）”。
% 入口未设置书名或作者时不写入。
\let\BookTitlePage@DefaultTitle\@title
\let\BookTitlePage@DefaultAuthor\@author
\AtBeginDocument{%
  \ifx\@title\BookTitlePage@DefaultTitle\else
    \ifx\BookTitlePage@Edition\@empty
      \hypersetup{pdftitle={\@title}}%
    \else
      \hypersetup{pdftitle={\@title（\BookTitlePage@Edition）}}%
    \fi
  \fi
  \ifx\@author\BookTitlePage@DefaultAuthor\else\hypersetup{pdfauthor={\@author}}\fi}
\makeatother

% 仓库地址：设置后，底部信息区首行为 GitHub 标志加地址，整行链接到仓库；提交哈希链接到 GitHub 上的
% 该提交“<仓库地址>/commit/<短哈希>”，外观不变（hidelinks）。地址不带末尾的 /，按原样读取，
% ~ _ % 等字符不需转义；不写或参数为空时没有地址行，哈希也不加链接。
\ExplSyntaxOn
\str_new:N \g__book_titlepage_repository_str
\NewDocumentCommand \RepositoryURL { v }
  { \str_gset:Nn \g__book_titlepage_repository_str {#1} }
\ExplSyntaxOff

% 提交信息文件由 scripts/commit-info.sh 经根目录 run.sh 写入编译输出目录，
% 四行依次为日期、短哈希、提交说明、clean 或 dirty。按原样逐行读取，
% 不经 TeX 转义；文件不存在或缺少日期、哈希时不显示提交信息。
\ExplSyntaxOn
\ior_new:N \g__book_titlepage_ior
\str_new:N \g__book_titlepage_date_str
\str_new:N \g__book_titlepage_hash_str
\str_new:N \g__book_titlepage_subject_str
\str_new:N \g__book_titlepage_state_str
\bool_new:N \g__book_titlepage_commit_bool
\cs_new_protected:Npn \__book_titlepage_get:N #1
  {
    \ior_if_eof:NTF \g__book_titlepage_ior
      { \str_gclear:N #1 }
      {
        \ior_str_get:NN \g__book_titlepage_ior \l_tmpa_str
        \str_gset_eq:NN #1 \l_tmpa_str
      }
  }
\NewDocumentCommand \TitlePageReadCommit { m }
  {
    \bool_gset_false:N \g__book_titlepage_commit_bool
    \file_if_exist:nT {#1}
      {
        \ior_open:Nn \g__book_titlepage_ior {#1}
        \__book_titlepage_get:N \g__book_titlepage_date_str
        \__book_titlepage_get:N \g__book_titlepage_hash_str
        \__book_titlepage_get:N \g__book_titlepage_subject_str
        \__book_titlepage_get:N \g__book_titlepage_state_str
        \ior_close:N \g__book_titlepage_ior
        \bool_lazy_or:nnF
          { \str_if_empty_p:N \g__book_titlepage_date_str }
          { \str_if_empty_p:N \g__book_titlepage_hash_str }
          { \bool_gset_true:N \g__book_titlepage_commit_bool }
      }
  }
% 默认读取本次编译的 <jobname>.commit；样张可在 \maketitle 前改读其他文件。
\AtBeginDocument { \TitlePageReadCommit { \c_sys_jobname_str .commit } }
% 有提交信息或设了仓库地址时，排底部信息区。
\NewDocumentCommand \BookTitlePageIfStampTF { +m +m }
  {
    \bool_lazy_or:nnTF
      { \g__book_titlepage_commit_bool }
      { ! \str_if_empty_p:N \g__book_titlepage_repository_str }
      {#1} {#2}
  }
% 仓库地址行（用户确认）：Nerd Fonts 的 GitHub 标志 oct-mark_github 加去掉开头 https:// 的地址，
% 与日期同为等宽次要灰；整行链接到仓库，标志与地址作为整体居中。标志取自等宽字体，字形宽 1 em，
% 步进却只有 0.6 em，所以放进 1 em 宽的盒子，再空 0.45 em 接地址。
\str_new:N \l__book_titlepage_repository_text_str
\cs_new_protected:Npn \__book_titlepage_repository_line:
  {
    \str_set_eq:NN \l__book_titlepage_repository_text_str \g__book_titlepage_repository_str
    \str_if_eq:eeT
      { \str_range:Nnn \l__book_titlepage_repository_text_str { 1 } { 8 } } { https:// }
      {
        \str_set:Ne \l__book_titlepage_repository_text_str
          { \str_range:Nnn \l__book_titlepage_repository_text_str { 9 } { -1 } }
      }
    \exp_args:Ne \href { \g__book_titlepage_repository_str }
      {
        \ttfamily
        \makebox [ 1em ] [ l ] { \symbol { "F408 } }
        \hspace { 0.45em }
        \str_use:N \l__book_titlepage_repository_text_str
      }
  }
% 底部信息区：设了仓库地址时首行为地址行；有提交信息时其下为眉题式三行：眉题“LAST COMMIT”；
% 日期与哈希，dirty 时附圆角细框标签；提交说明。不用分隔符串联，依靠换行、留白及字体差异区分标签与数值。
% 只有哈希本身链接到提交；dirty 时同样链接到所基于的提交。
% 间距（用户确认）：行高 11 pt；地址行到眉题的基线间距为 1.5 倍行高，眉题与它标注的两行按行高排成一组，
% 组间空白约为组内的 2 倍。dirty 时标签细框高出数字约 1 pt、低于基线约 2 pt，日期行上下各补足这么多，
% 使细框上下的空白与没有标签时接近。
\bool_new:N \l__book_titlepage_dirty_bool
\NewDocumentCommand \BookTitlePageStampText { }
  {
    \str_if_empty:NF \g__book_titlepage_repository_str
      {
        \__book_titlepage_repository_line:
        \bool_if:NT \g__book_titlepage_commit_bool { \\ [ 5.5pt ] }
      }
    \bool_if:NT \g__book_titlepage_commit_bool
      {
        \str_if_eq:VnTF \g__book_titlepage_state_str { dirty }
          { \bool_set_true:N \l__book_titlepage_dirty_bool }
          { \bool_set_false:N \l__book_titlepage_dirty_bool }
        { \BookTitlePageEyebrow \fontsize { 6.5 } { 8 } \selectfont LAST ~ COMMIT }
        \bool_if:NTF \l__book_titlepage_dirty_bool { \\ [ 1pt ] } { \\ }
        { \ttfamily \str_use:N \g__book_titlepage_date_str }
        \hspace { 0.8em }
        { \ttfamily \color { BookTitlePageAccent }
          \str_if_empty:NTF \g__book_titlepage_repository_str
            { \str_use:N \g__book_titlepage_hash_str }
            {
              \exp_args:Ne \href
                { \g__book_titlepage_repository_str / commit / \g__book_titlepage_hash_str }
                { \str_use:N \g__book_titlepage_hash_str }
            }
        }
        \bool_if:NT \l__book_titlepage_dirty_bool
          { \hspace { 0.6em } \usebox \BookTitlePageDirtyBox }
        \str_if_empty:NF \g__book_titlepage_subject_str
          {
            \bool_if:NTF \l__book_titlepage_dirty_bool { \\ [ 2pt ] } { \\ }
            \BookTitlePagePlainLatin \str_use:N \g__book_titlepage_subject_str
          }
      }
  }
\ExplSyntaxOff
% 未提交修改标签：0.4 pt 浅灰圆角细框，与代码卡片同一视觉语言。
% 在扉页的 tikzpicture 之外预先排入盒子，避免嵌套 TikZ。
\newsavebox{\BookTitlePageDirtyBox}
\newcommand*{\BookTitlePageMakeDirtyBox}{%
  \sbox{\BookTitlePageDirtyBox}{\tikz[baseline=(tag.base)]
    \node[draw=BookTitlePageHairline,line width=0.4pt,rounded corners=1.1mm,
      inner xsep=1.3mm,inner ysep=0.45mm,text=BookTitlePageInk,
      font=\sffamily\fontsize{6.5}{8}\selectfont] (tag) {含未提交修改};}}

% 版面以纸张中心为原点：标题区在上，配图居中偏下，底部信息区贴近底边；双栏正文不影响绝对定位的扉页。
% 版本标记（用户确认，2026-10-15）：全角括号、苹方常规体 22 pt（书名字号的一半）、
% 次要灰，与书名同基线，悬挂在书名右侧、不参与居中；书名仍以页面中轴居中，
% 与英文书名、作者对齐。
\makeatletter
\renewcommand{\maketitle}{%
  \clearpage
  \thispagestyle{empty}%
  \BookTitlePageMakeDirtyBox
  \null
  \begin{tikzpicture}[remember picture,overlay]
    \begin{scope}[shift={(current page.center)}]
      \node[anchor=base,inner sep=0pt,text=BookTitlePageInk,
        font=\sffamily\bfseries\fontsize{44}{52}\selectfont] (title) at (0,\BookTitlePageTitleY) {\@title};
      \ifx\BookTitlePage@Edition\@empty\else
        \node[anchor=base west,inner sep=0pt,text=BookTitlePageSecondary,
          font=\sffamily\mdseries\fontsize{22}{26}\selectfont] at (title.base east)
          {（\BookTitlePage@Edition）};
      \fi
      \ifx\BookTitlePage@EnglishTitle\@empty\else
        \node[anchor=base,inner sep=0pt,text=BookTitlePageSecondary,
          font=\sffamily\fontsize{15}{18}\selectfont] at (0,\BookTitlePageEnglishY)
          {\BookTitlePage@EnglishTitle};
      \fi
      \node[anchor=base,inner sep=0pt,text=BookTitlePageInk,
        font=\BookTitlePageMediumLatin\BookTitlePageMediumCJK\fontsize{12}{15}\selectfont]
        at (0,\BookTitlePageAuthorY) {\@author};
      \ifx\BookTitlePage@Image\@empty\else
        \node[inner sep=0pt] at (0,\BookTitlePageImageY) {\includegraphics[
          width=\BookTitlePageImageWidth,height=\BookTitlePageImageHeight,
          keepaspectratio]{\BookTitlePage@Image}};
      \fi
      % 底边对齐：整块从底边向上排，地址行在提交信息之上，提交说明换行也向上增长，不进入下边距。
      % flush center 使用 fil 胶居中；align=center 的有限伸缩会拉开 xeCJK 字距。
      \BookTitlePageIfStampTF{%
        \node[anchor=south,inner sep=0pt,text width=150mm,align=flush center,
          text=BookTitlePageSecondary,font=\sffamily\fontsize{8}{11}\selectfont]
          at (0,\BookTitlePageStampY) {\BookTitlePageStampText};
      }{}%
    \end{scope}
  \end{tikzpicture}%
  \clearpage}
\makeatother
